
\subsection{Optimized Peeling for $r \ge 3$}
\label{sec:opt_r3}

For $r \ge 3$, path splitting (Algorithm~\ref{alg:removeRClique}) is the computational bottleneck:
it branches over each pivot in the conflict set, yielding worst-case exponential time.
We introduce \emph{dead-path detection}---a polynomial-time test that identifies paths where no valid subpath of size $\ge s$ can survive.
When the test succeeds, we skip the exponential splitting entirely and subtract the support directly.

\stitle{Key Idea: Budgeted Hitting Set.}
Path splitting is equivalent to a \emph{budgeted hitting set} problem.
Given a set of removed $r$-cliques $C_{\mathrm{rm}} = \{C_0, \dots, C_{m-1}\}$ on path $\TPath$,
each $C_i$ induces a \emph{conflict set} $F_i = C_i \cap V_p(\TPath)$:
the pivot members that must be ``hit'' (at least one removed) to invalidate $C_i$.
The \emph{budget} is $B = |V_p(\TPath)| - (s - |V_h(\TPath)|)$: the maximum number of pivots that may be removed while the path still encodes at least one $s$-clique.

A path is \emph{dead} if and only if $\tau(\mathcal{F}) > B$, where $\tau(\mathcal{F})$ is the minimum hitting set size of the conflict family $\mathcal{F} = \{F_0, \dots, F_{m-1}\}$.
Computing $\tau$ exactly is NP-hard, but the following polynomial-time lower bounds suffice for the vast majority of practical instances.

\begin{lemma}[Forced Vertex]\label{lem:forced_v}
If a pivot vertex $v$ appears in every $r$-clique on $\TPath$ (i.e., its conflict degree equals $\binom{|\TPath|-1}{r-1}$), then $v$ must be removed in every feasible solution.
If $v$ is a hold vertex, the path is unconditionally dead.
\end{lemma}

\begin{proof}
Since $v$ belongs to every $r$-clique, every conflict set $F_i$ contains $v$.
The only way to hit all $F_i$ is to remove $v$.
If $v$ is a hold, removing it kills the path (Theorem~\ref{lem:path_vertex_remove}).
\end{proof}

\begin{lemma}[Singleton Propagation]\label{lem:singleton_v}
If $|F_i| = 1$, say $F_i = \{p\}$, then $p$ must be removed in every feasible solution.
Removing $p$ reduces the budget by one and may create new singletons in the remaining conflict sets.
Iterative propagation runs in $O(m \cdot r)$ time.
If at any point $|V_p| + |V_h| < s$, the path is dead.
\end{lemma}

\begin{proof}
A singleton conflict set has exactly one pivot.
Any hitting set must include that pivot.
After removing it, any $F_j$ containing $p$ shrinks by one, potentially becoming a new singleton.
\end{proof}

\begin{lemma}[Disjoint Packing Lower Bound]\label{lem:packing_v}
If there exist $t$ pairwise disjoint sets in $\mathcal{F}$, then any feasible hitting set has size $\ge t$.
If $t > B$, the path is dead.
A greedy algorithm finds $t$ in $O(m^2)$ time.
\end{lemma}

\begin{proof}
Each disjoint conflict set requires at least one distinct pivot to be removed.
Hence $\tau(\mathcal{F}) \ge t$.
\end{proof}

\begin{lemma}[Component Decomposition]\label{lem:component_v}
Let $G_c$ be the \emph{conflict graph} on $V_p$, where two pivots are adjacent iff they co-occur in some $F_i$.
If $G_c$ has connected components $\mathcal{C}_1, \dots, \mathcal{C}_k$, then:
\[
\tau(\mathcal{F}) = \sum_{j=1}^{k} \tau(\mathcal{F}[\mathcal{C}_j]),
\]
where $\mathcal{F}[\mathcal{C}_j]$ denotes the conflict sets restricted to component $\mathcal{C}_j$.
Computing the packing lower bound per component yields a bound $\ge$ the global packing bound.
\end{lemma}

\begin{proof}
Since no conflict set spans two components, the hitting set decomposes independently.
Per-component packing sums to at least the global packing (each disjoint pair is within one component).
\end{proof}


\stitle{Algorithm.}
These conditions are applied as a preprocessing step before the recursive path split (Algorithm~\ref{alg:removeRClique}).
Algorithm~\ref{alg:pathsplit_pruning} summarizes the pruned procedure:

\begin{enumerate}[leftmargin=*]
\item \textbf{Forced vertex detection}: remove pivots with conflict degree $= \binom{|\TPath|-1}{r-1}$.
\item \textbf{Singleton propagation}: iteratively remove singleton conflict pivots; propagate until stable.
\item \textbf{Packing + component bound}: compute per-component disjoint packing; if total $> B$, declare dead.
\item If dead: subtract support for all $r$-cliques on $\TPath$ via direct enumeration in $O\bigl(\binom{|\TPath|}{r}\bigr)$.
\item If alive: proceed with recursive path split (Algorithm~\ref{alg:removeRClique}).
\end{enumerate}


\begin{algorithm}[t]
\small
\caption{\textsc{PrunedUpdateIndex}}
\label{alg:pathsplit_pruning}
\KwIn{clique path $\TPath{} = \{V_H, V_P\}$, $r$-cliques to remove $C_{\mathrm{rm}}$, clique size $s$}
\KwOut{All subpaths after deletion, or $\emptyset$ if dead}
\SetKwFunction{Split}{PathSplit}
\SetKwProg{Fn}{Function}{}{}

$F_i \gets C_i \cap V_P$ for each $C_i \in C_{\mathrm{rm}}$\tcp*{conflict sets}
$B \gets |V_P| - (s - |V_H|)$\tcp*{removal budget}

\tcp{Forced vertex (Lemma~\ref{lem:forced_vertex})}
\ForEach{$v$ with conflict degree $\ge \binom{|\TPath|-1}{r-1}$}{
  \lIf{$v \in V_H$}{\Return $\emptyset$}
  $V_P \gets V_P \setminus \{v\}$; resolve all $F_i \ni v$\;
}

\tcp{Singleton propagation (Lemma~\ref{lem:singleton})}
\Repeat{no change}{
  \ForEach{active $F_i$ with $|F_i| = 1$, say $\{p\}$}{
    $V_P \gets V_P \setminus \{p\}$; resolve all $F_j \ni p$\;
    \lIf{$|V_H| + |V_P| < s$}{\Return $\emptyset$}
  }
}

\tcp{Component packing LB (Lemmas~\ref{lem:packing},~\ref{lem:component_lb})}
Build conflict graph $G_c$; find components $C_1, \dots, C_k$\;
$\mathit{LB} \gets \sum_{j=1}^{k} \textsc{GreedyPacking}(\{F_i \mid F_i \subseteq C_j\})$\;
\lIf{$\mathit{LB} > B$}{\Return $\emptyset$}

\Return \Split{$V_H, V_P$} with reduced $\mathcal{F}$\tcp*{Alg.~\ref{alg:removeRClique}}
\end{algorithm}



\stitle{Effectiveness.}
Table~\ref{tab:leafdeath} reports the fraction of paths detected as dead by the above tests.
On representative datasets at $s = 20$, dead-path detection prunes 95--97\% of all paths,
reducing the number of expensive recursive splits by up to $47\times$.
This is the primary reason our method completes on parameter ranges where the baseline times out.

\stitle{Complexity.}
The preprocessing (Lemmas~\ref{lem:forced_v}--\ref{lem:component_v}) runs in $O(m^2 \cdot r)$ time, where $m = |C_{\mathrm{rm}}|$.
When the path is dead, the support subtraction costs $O\bigl(\binom{|\TPath|}{r}\bigr)$---polynomial in the path size.
The exponential recursive split (Algorithm~\ref{alg:removeRClique}) is invoked only for the surviving 3--5\% of paths.
